% 0909 收尾轮：\zhenti 独占行缺陷最小复现＋变体实验（列宽 82.8mm＝工程同参）
\documentclass{article}
\usepackage[margin=17.575mm]{geometry}
\usepackage{fontspec}
\usepackage{xeCJK}
\usepackage{multicol}
\setCJKmainfont[Path=C:/Users/28120/AppData/Local/Microsoft/Windows/Fonts/]{FZShuSong-Z01S.ttf}
\setCJKfamilyfont{fangsong}[Path=fonts/]{FZFSK.TTF}
\newfontfamily{\fangsonglat}[Path=fonts/]{FZFSK.TTF}
\xeCJKsetup{CheckSingle=true,PunctStyle=plain}
\newcommand{\ansul}[1]{#1}
\newcommand{\cha}{{\fangsonglat\symbol{"00D7}}}
\newcommand{\gou}{{\fangsonglat\symbol{"221A}}}
\newlength{\colw}
\setlength{\colw}{82.8mm}
% V1＝现工程形
\newcommand{\ztA}[2]{{\parfillskip=0pt\emergencystretch=3em\noindent#1\nobreak\mbox{}\nobreak\hspace{0pt plus 1fil}(\makebox[1.8em]{\ansul{#2}})\hspace{0.56mm}\par}}
% V2＝mbox 去掉、nobreak 挪到 fil 之后
\newcommand{\ztB}[2]{{\parfillskip=0pt\emergencystretch=3em\noindent#1\nobreak\hspace{0pt plus 1fil}\nobreak(\makebox[1.8em]{\ansul{#2}})\hspace{0.56mm}\par}}
% V3＝括号组整体 \mbox 胶囊（fil 在外），题干高容差
\newcommand{\ztC}[2]{{\parfillskip=0pt\emergencystretch=3em\tolerance=9999\noindent#1\nobreak\hspace{0pt plus 1fil}\mbox{(\makebox[1.8em]{\ansul{#2}})\hspace{0.56mm}}\par}}
% V4＝\hfil 换 \hspace fil＋末尾 \parfillskip 恒等、去 mbox、题干尾 nobreak 直连 fil，容差放开＋pretolerance 关
\newcommand{\ztD}[2]{{\parfillskip=0pt\emergencystretch=5em\tolerance=9999\pretolerance=-1\noindent#1\nobreak\hspace{0pt plus 1fil}\nobreak\mbox{(\makebox[1.8em]{\ansul{#2}})}\hspace{0.56mm}\par}}
\begin{document}
\fontsize{10.09pt}{18.25pt}\selectfont
\begin{multicols}{2}
\noindent A1\quad\ztA{两个空间向量的模相等，则这两个向量相等．}{\cha{}}
\noindent A2\quad\ztA{相等向量一定是共线向量．}{\gou{}}
\noindent A3\quad\ztA{若\(\overrightarrow{a}\parallel\overrightarrow{b}\)，\(\overrightarrow{b}\parallel\overrightarrow{c}\)，则\nobreak\(\overrightarrow{a}\parallel\overrightarrow{c}\)．}{\cha{}}
\noindent A4\quad\ztA{若\(\overrightarrow{a}\parallel\overrightarrow{b}\)，则存在唯一实数\(\lambda\)，使\(\overrightarrow{a}=\lambda\overrightarrow{b}\)．}{\cha{}}
\noindent A5\quad\ztA{若\(\overrightarrow{p}=x\overrightarrow{a}+y\overrightarrow{b}\)(\(\overrightarrow{a}\)，\(\overrightarrow{b}\)不共线)，则\nobreak\(\overrightarrow{p}\)，\(\overrightarrow{a}\)，\(\overrightarrow{b}\)共面．}{\gou{}}
\end{multicols}
\begin{multicols}{2}
\noindent B1\quad\ztB{两个空间向量的模相等，则这两个向量相等．}{\cha{}}
\noindent B2\quad\ztB{相等向量一定是共线向量．}{\gou{}}
\noindent B3\quad\ztB{若\(\overrightarrow{a}\parallel\overrightarrow{b}\)，\(\overrightarrow{b}\parallel\overrightarrow{c}\)，则\nobreak\(\overrightarrow{a}\parallel\overrightarrow{c}\)．}{\cha{}}
\noindent B4\quad\ztB{若\(\overrightarrow{a}\parallel\overrightarrow{b}\)，则存在唯一实数\(\lambda\)，使\(\overrightarrow{a}=\lambda\overrightarrow{b}\)．}{\cha{}}
\noindent B5\quad\ztB{若\(\overrightarrow{p}=x\overrightarrow{a}+y\overrightarrow{b}\)(\(\overrightarrow{a}\)，\(\overrightarrow{b}\)不共线)，则\nobreak\(\overrightarrow{p}\)，\(\overrightarrow{a}\)，\(\overrightarrow{b}\)共面．}{\gou{}}
\end{multicols}
\begin{multicols}{2}
\noindent C1\quad\ztC{两个空间向量的模相等，则这两个向量相等．}{\cha{}}
\noindent C2\quad\ztC{相等向量一定是共线向量．}{\gou{}}
\noindent C3\quad\ztC{若\(\overrightarrow{a}\parallel\overrightarrow{b}\)，\(\overrightarrow{b}\parallel\overrightarrow{c}\)，则\nobreak\(\overrightarrow{a}\parallel\overrightarrow{c}\)．}{\cha{}}
\noindent C4\quad\ztC{若\(\overrightarrow{a}\parallel\overrightarrow{b}\)，则存在唯一实数\(\lambda\)，使\(\overrightarrow{a}=\lambda\overrightarrow{b}\)．}{\cha{}}
\noindent C5\quad\ztC{若\(\overrightarrow{p}=x\overrightarrow{a}+y\overrightarrow{b}\)(\(\overrightarrow{a}\)，\(\overrightarrow{b}\)不共线)，则\nobreak\(\overrightarrow{p}\)，\(\overrightarrow{a}\)，\(\overrightarrow{b}\)共面．}{\gou{}}
\end{multicols}
\begin{multicols}{2}
\noindent D1\quad\ztD{两个空间向量的模相等，则这两个向量相等．}{\cha{}}
\noindent D2\quad\ztD{相等向量一定是共线向量．}{\gou{}}
\noindent D3\quad\ztD{若\(\overrightarrow{a}\parallel\overrightarrow{b}\)，\(\overrightarrow{b}\parallel\overrightarrow{c}\)，则\nobreak\(\overrightarrow{a}\parallel\overrightarrow{c}\)．}{\cha{}}
\noindent D4\quad\ztD{若\(\overrightarrow{a}\parallel\overrightarrow{b}\)，则存在唯一实数\(\lambda\)，使\(\overrightarrow{a}=\lambda\overrightarrow{b}\)．}{\cha{}}
\noindent D5\quad\ztD{若\(\overrightarrow{p}=x\overrightarrow{a}+y\overrightarrow{b}\)(\(\overrightarrow{a}\)，\(\overrightarrow{b}\)不共线)，则\nobreak\(\overrightarrow{p}\)，\(\overrightarrow{a}\)，\(\overrightarrow{b}\)共面．}{\gou{}}
\end{multicols}
\end{document}
